\begin{theorem}[指数温度调度下的 Perron 自由能相变律]\label{thm:pom-replica-softcore-temperature-free-energy-phase-transition}
\leanverified{paper\_pom\_replica\_softcore\_temperature\_free\_energy\_phase\_transition}
令 \(\varphi=\frac{1+\sqrt5}{2}\)，并令 \(\rho(\cdot)\) 表示谱半径。
对任意常数 \(\alpha\ge 0\)，取温度参数随阶数指数衰减
\[
p_q:=2^{-\alpha q}\qquad(q\in\NN).
\]
则极限
\[
F(\alpha):=\lim_{q\to\infty}\frac1q\log \rho\!\bigl(T_{p_q}^{(q)}\bigr)
\]
存在，且满足闭式
\[
F(\alpha)=\max\bigl\{\log\varphi,\ (1-\alpha)\log 2\bigr\}.
\]
等价地，存在临界指数
\[
\alpha_c:=1-\log_2\varphi
\]
使得当 \(\alpha<\alpha_c\) 时 \(F(\alpha)=(1-\alpha)\log 2\)，当 \(\alpha\ge\alpha_c\) 时 \(F(\alpha)=\log\varphi\)。
\end{theorem}

\begin{proof}
由 \eqref{eq:pom-replica-softcore-temperature-decomposition}，
\(
T^{(q)}_{p}=p\,\mathbf1\mathbf1^{\mathsf T}+(1-p)K^{\otimes q}
\)。
注意到 \(\mathbf1\mathbf1^{\mathsf T}\) 为秩一非负矩阵，其 Perron 根为 \(\mathbf1^{\mathsf T}\mathbf1=2^q\)，故
\[
\rho\!\bigl(p\,\mathbf1\mathbf1^{\mathsf T}\bigr)=p\,2^q.
\]
又由于张量谱乘法性与 \(\rho(K)=\varphi\)，有
\[
\rho(K^{\otimes q})=\rho(K)^q=\varphi^q,
\qquad
\rho\!\bigl((1-p)K^{\otimes q}\bigr)=(1-p)\varphi^q.
\]
并且按元素非负性有 \(T_p^{(q)}\ge p\,\mathbf1\mathbf1^{\mathsf T}\) 与 \(T_p^{(q)}\ge(1-p)K^{\otimes q}\)，故由谱半径对非负矩阵的单调性得到下界
\[
\rho(T_p^{(q)})\ \ge\ \max\bigl\{p\,2^q,\ (1-p)\varphi^q\bigr\}.
\]
另一方面，\(T_p^{(q)}\) 为对称非负矩阵，故 \(\rho(T_p^{(q)})=\|T_p^{(q)}\|_2\)。
由算子范数三角不等式得
\[
\rho(T_p^{(q)})
\le
\bigl\|p\,\mathbf1\mathbf1^{\mathsf T}\bigr\|_2+\bigl\|(1-p)K^{\otimes q}\bigr\|_2
=p\,2^q+(1-p)\varphi^q.
\]
因此对任意 \(p\in[0,1]\) 有夹逼
\[
\max\bigl\{p\,2^q,\ (1-p)\varphi^q\bigr\}
\ \le\
\rho(T_p^{(q)})
\ \le\
p\,2^q+(1-p)\varphi^q.
\]
取 \(p=p_q=2^{-\alpha q}\) 得 \(p_q\,2^q=2^{(1-\alpha)q}\)，且 \((1-p_q)\varphi^q=\varphi^q(1+o(1))\)。
由两项指数增长的最大指数主导律，
\[
\lim_{q\to\infty}\frac1q\log\bigl(2^{(1-\alpha)q}+\varphi^q\bigr)
=\max\bigl\{(1-\alpha)\log 2,\ \log\varphi\bigr\},
\]
结合夹逼即得所示闭式。
临界指数由 \((1-\alpha)\log2=\log\varphi\) 等价变形得到。证毕。
\end{proof}
